\section{考察}\label{sec:discussion}

\subsection{種の価値の差はいつ現れるか}

同じカット値の候補の価値の差（C1）は，CIM の出力が多様なすべての問題で，種の近くを探す SA（$T_0=0.5$）では予算によらず成り立った．一方，開始温度が高い SA や長い予算では差が消えた．事前登録の実行で測った SA の出力と種の距離は，差が残った G22 で頂点の 26\%，差が消えた G55・G70 で 43〜44\% だった．全反転を同一視した無作為な 2 状態の距離の期待値は約 50\% なので，差が消えた条件では SA は種をほとんど忘れている．種の価値は，後段の探索が種の情報をどれだけ保つかで決まる量であり，候補だけの性質ではない．

小規模問題の全列挙では，価値は最適解までの距離（経路長）でほぼ決まり，障壁は弱かった．大規模問題でも，SA が種の近くを探す条件では「最適解への近さ」が効いていると考えるのが自然である．ただし，大規模問題では最適解が分からないので，この解釈は直接は確かめられない．

\subsection{固有ベクトル指標はなぜ一部の問題で効くか}

$E$ は G22 で事前登録どおりに再現し，G1 でも同じ大きさで成り立ったが，G23 や小規模問題の $\pm1$ 重みのグラフでは効かなかった．CIM の増幅の初期には $A$ の最小固有値側の固有ベクトルが最も速く成長し，ポンプを速く立ち上げて線形増幅の段階を短くすると，$u_1$ の系統に入る候補の割合が下がった（\ref{sec:ga} 節）．したがって $E$ は，候補が増幅の初期にどのモードに引き込まれたかの痕跡である．小規模問題で $a_1$ が最適解までの距離と相関したことから，$u_1$ の系統に入ることは，問題によっては最適解に近い領域に入ることを意味する．どの問題でそうなるかは，最適解と $u_1$ の符号パターンの重なりに依存するはずだが，本論文ではこれを予測する方法は見つけていない．

局所的な指標の向きが小規模問題と G-set で逆になったことも，指標の意味が問題の構造に依存することを示している．重み 1 のランダムグラフでは利得 $-2$〜$0$ の頂点が多いことが「動ける余地」として正に効き，小規模問題では「未完成の局所構造」として負に効いた可能性があるが，これは事後的な解釈である．問題をまたいで同じ向きに効いたのは試し打ちだけだった．試し打ちは後段の探索そのものの縮小版なので，問題の構造によらず効くのは自然である．

\subsection{なぜ選別は役に立たないか}

C1 と C2 が成り立つ問題でも，同じ総時間で比べると，カット値を超える選別の効果はほとんどなかった．理由は 3 つ考えられる．

第一に，カット値で説明できない価値の差は，カット値による差より小さい．無作為な選び方とカット値による選び方の差（主条件で BKS との差の平均 0.08\%）に比べ，カット値を超える選別で動かせた差（最大 0.003\%）は一桁以上小さかった．

第二に，評価量は $m$ 本の SA の\textbf{最良値}であり，価値（期待値）の差がそのまま最良値の差にはならない．カット値の上位の候補はすでに価値の高い群に入っていて，その中での順位の入れ替えは最良値をほとんど変えない．

第三に，選別の費用がかかる．LS と局所的な指標は SA に比べて十分安いが，試し打ちは短い予算では SA の予算全体より高くつく．固有ベクトルの計算は問題ごとに 1 回で済むが，それでも小さな $m$ では SA の本数を削る．

\subsection{CIM の種はいつ役に立つか}

CIM の種とランダム初期化の優劣は，後段の予算で逆転した．後段の予算が短いと，SA はランダムな初期解から良い解まで進めず，CIM の出力を出発点にできることの利点が大きい（G70 の 500 掃引では BKS との差で 0.79\% の改善）．後段の予算が長いと，SA はランダムな初期解からでも同程度の解に達し，CIM の生成時間を SA の本数に回したほうが最良値が伸びる．メメティックアルゴリズムでは，同じ世代数で比べると CIM の初期集団が G55・G70 で良かった（\ref{sec:ga} 節）．GA の後段は SA より強い探索だが，世代数が限られる場合には CIM の初期集団の品質が効くと考えられる．ただし，この比較は CIM の生成時間を含めていない．

実用上の示唆は次のとおりである．CIM を後段の探索の初期解生成に使うなら，(1) 候補はカット値（LS 後のカット値）で選べば十分であり，それ以上の選別に時間を使う価値は小さい．(2) CIM の生成時間を後段に回した場合との比較を必ず行うべきであり，後段の予算が長い設定では CIM を使わないほうが良いことがある．

\subsection{限界}

\begin{itemize}
  \item 小規模問題では，CIM がほぼ常に最適解を出すので，増幅の途中で止めた CIM の出力を候補とした．大規模問題の CIM の出力と性質が同じとは限らない．
  \item SA は途中の最良状態（初期解を含む）を返すので，開始温度が高く予算が短いとラベルが初期カット値と一致する．この条件は判定不能とし，判定不能の基準は結果を見てから改めた（\ref{sec:prereg}，\ref{sec:gset} 節）．
  \item 時間はすべて 1 台の計算機（24 スレッド）での実測であり，CIM は物理的な装置ではなく数値シミュレーションである．物理的な CIM では生成時間がはるかに短い可能性があり，その場合はランダム初期化との比較の結論が変わりうる．
  \item 平面グラフの G15・G39 では CIM の出力が 1〜5 状態に集まり，評価から除いた．CIM のパラメータ（結合の強さ以外）は G22 で調整した値であり，問題ごとに調整すれば多様性が増す可能性がある．
  \item G77 は計算時間の都合で除いた．比較は 7 問題で行い，問題を単位にした検定の検出力は低い．
  \item BKS は以前の研究で用いた既知の最良値の表に従った．値の出典によって BKS との差の絶対値は変わりうるが，同じ問題の中での方法の比較には影響しない．
\end{itemize}

\subsection{結論}

同じカット値の CIM 解は，後段の SA で同じ価値を持たない．この差は，SA が種の情報を保つ設定で，CIM の出力が多様なすべての問題で再現した．差の一部は，CIM の増幅の初期にどの固有ベクトルの系統に入ったかや，短い SA の試し打ちで予測できたが，効く指標は問題によって違った．そして，CIM の生成を含めた同じ総時間で比べると，カット値を超える選別は到達品質をほとんど改善しなかった．CIM の種が役に立つかどうかは，選別よりも，後段の探索にどれだけの予算を使うかで決まる．
